% =====================================================================
%  Supplementary appendix (TEXT-REVISER, revision 1): material moved out of
%  Sections 6 and 7 to shorten the main text. Every label is preserved.
% =====================================================================
\section{Supplement to Section~\ref{sec:estimation}}\label{app:est-supp}

This appendix collects the material behind \cref{sec:estimation} that the main text cites:
the treatment of the observation window in EM, the Gibbs sampler, nonparametric kernels,
the classical boundary and time-rescaling results with their conditions, the Davies
problem, and the proofs of the results of \cref{sec:estimation}.

\subsection{Edge effects}\label{app:est-edge}

\begin{remark}[Edge effects]\label{rem:est-edge}
The finite window enters the M-step only through the window terms $\mathcal G_{kj}$ and
$\Psi^T_k(e)$. Ignoring them means setting $\mathcal G_{kj}=n_j$ and $\Psi^T_k(e)=1$, and
letting $m_j(\beta)=1/\beta$. This gives the familiar updates $\alpha_{kj}=N_{kj}/n_j$,
$\beta_{kj}=N_{kj}/W_{kj}$ and $\beta_k=\sum_eN_{ke}/\sum_eS_e$. The resulting iteration is
an EM algorithm for an edge-free pseudo-likelihood, in which every incident's offspring and
every shock's response are counted in full, and it is monotone for that criterion. It is
not an EM algorithm for $\ell_T$: its $\ell_T$-trace need not increase, and its fixed point
is biased. Offspring of incidents near $T$, and the responses to shocks near $T$, fall
partly after $T$. As a result, $\alpha_{kj}$ and $\beta_k$ are underestimated, and decay
rates are overestimated because the observed lags are right-truncated. The relative error
in $\alpha_{kj}$ is about $n_j^{-1}\sum_le^{-\beta_{kj}(T-t^j_l)}\approx1/(\beta_{kj}T)$. It
is negligible when the kernel is short relative to the window. It is not negligible for a
salient event shortly before $T$, which is exactly the case of real-time attribution
(\cref{sec:attribution,sec:forecasting}). For ETAS models, a closely related approximation
replaces the integral of the triggering function over the observation region by its
integral over the whole domain; \citet{schoenberg2013} discusses it. The exact window terms
cost nothing extra here, and we use them throughout.

Incidents before time $0$ excite the start of the window. We therefore condition on a
burn-in segment $[0,T_0]$, whose incidents enter the intensities as history but contribute
no likelihood terms. The window terms then become
\[
  \mathcal G_{kj}=\sum_{l}\int_{(T_0-t^j_l)^+}^{T-t^j_l}g_{kj}(u)\,\dd u,\qquad
  \Psi^T_k(e)=\Psi_k(T-\tau_e)-\Psi_k(T_0-\tau_e),
\]
with $\int_{T_0}^T\mu_k$ in place of $\int_0^T\mu_k$, and the E-step sums run over
incidents in $(T_0,T]$. \Cref{prop:est-mstep} holds unchanged, with
$\nu(\dd u)=\sum_l\ind{(T_0-t^j_l)^+\le u\le T-t^j_l}\,\dd u$ in its proof. This $\nu$ is
again a positive measure that is not a point mass, so its Laplace transform is strictly
log-convex. The function $m_j$ is the mean of $u$ under $e^{-\beta u}\nu(\dd u)$, and
$m_j(0+)$ is the mean of $\nu/\nu(\mathbb R)$.
\end{remark}

\subsection{Gibbs sampling by data augmentation}\label{app:est-gibbs}

Given the labels $Z$, the
complete-data likelihood \eqref{eq:est-complete} is a product of Poisson-process
likelihoods. Each parameter then has a simple full conditional. One sweep consists of the
following steps.
\begin{enumerate}[(1)]
  \item \emph{Labels.} Draw each $Z^k_i$ from the E-step probabilities: first the category
    (background, shock $e$, or stream $j$), in $O(M+E)$ operations; then, given stream
    $j$, the parent $l$ with probability proportional to $e^{-\beta_{kj}(t^k_i-t^j_l)}$.
    The second draw is a binary search on rescaled cumulative sums, in $O(\log n_{\mathrm{ev}})$
    operations. Record the realised counts $N_{kj}$, $N_{ke}$, $N^0_k$ and lag sums
    $W_{kj}$, $W_{ke}$.
  \item \emph{Branching ratios.} The complete-data likelihood in $\alpha_{kj}$ is
    $\alpha_{kj}^{N_{kj}}e^{-\alpha_{kj}G_{kj}}$, and this includes the window term
    exactly. Hence
    \[
      \alpha_{kj}\mid\cdot\ \sim\ \operatorname{Gamma}\big(a_1+N_{kj},\ \omega_{kj}+G_{kj}\big)
      \ \text{truncated to}\ \big[0,\,1/\Mmat^{(-kj)}_{jk}\big),
    \]
    which is sampled by inversion (\cref{lem:est-trunc}). The scales follow from
    $\omega_{kj}\mid\cdot\sim\operatorname{Gamma}(a_1+a_2,\alpha_{kj}+\varsigma)$ and
    $\varsigma\mid\cdot\sim\operatorname{Gamma}(a_\varsigma+a_2|\mathcal P|,\,b_\varsigma+\sum\omega_{kj})$,
    where $\mathcal P$ is the set of shrunk entries. (In the spike-and-slab version, the
    slab replaces $(a_1,\omega_{kj})$ by $(a_{\mathrm s},b_{\mathrm s})$.)
  \item \emph{Spike and slab.} Update $(\iota_{kj},Z^k)$ as a block. First draw
    $\iota_{kj}$ with $Z^k$ integrated out. The odds are
    $\pi_{kj}\exp\{\ell_{T,k}(\tilde\alpha_{kj})\}/[(1-\pi_{kj})\exp\{\ell_{T,k}(0)\}]$, where
    $\ell_{T,k}(x)$ is the stream-$k$ log-likelihood with $\alpha_{kj}=x$, computable in
    $O(n_{\mathrm{ev}})$ by \eqref{eq:est-recursion}.
    Set $\iota_{kj}=1$ only if the truncation allows it. Then draw $Z^k$ given
    $\iota_{kj}$. When $\iota_{kj}=0$, draw $\tilde\alpha_{kj}$ from its prior. Updating $\iota_{kj}$ given $Z^k$ instead would mix badly: any label
    pointing to stream $j$ forces $\iota_{kj}=1$.
  \item \emph{Loadings and background.} With gamma priors,
    $\beta_k\mid\cdot\sim\operatorname{Gamma}(a_\beta+\sum_eN_{ke},\,b_\beta+\sum_eS_e\Psi^T_k(e))$.
    A constant background has $\mu_k\mid\cdot\sim\operatorname{Gamma}(a_\mu+N^0_k,\,b_\mu+T)$.
    For the spline background, a Gaussian prior $\theta_k\sim N(0,(\lambda^{\mathrm{pen}}_kP+\varepsilon\Imat)^{-1})$
    is the Bayesian form of \eqref{eq:est-pen}. Its full conditional, that of a Poisson
    process with log-linear intensity, is log-concave and is updated by a Langevin or
    Laplace-proposal Metropolis step.
  \item \emph{Decay rates.} The full conditional density of $\beta_{kj}$ is proportional
    to
    \[
      \pi(\beta)\,\beta^{N_{kj}}e^{-\beta W_{kj}}\exp\Big\{\alpha_{kj}\sum_le^{-\beta(T-t^j_l)}\Big\}.
    \]
    The last factor is the window correction. We update $\log\beta_{kj}$ by slice sampling
    \citep{neal2003}. When the window correction is negligible and $\pi$ is a gamma
    density, the update is conjugate.
\end{enumerate}
Each posterior draw of $\vartheta$ gives draws of $\spr(\Amat)$, $\Mmat$ and the surge
decomposition of \cref{thm:attribution}, so parameter uncertainty propagates to every
reported quantity (\cref{sec:at-uncertainty}).

\subsection{Nonparametric kernels and sparse excitation}\label{sec:est-nonpar}

\paragraph{Histogram and spline kernels by EM.} Let $\varphi_{kj}(u)=\sum_{m=1}^{L}c_{kj,m}b_m(u)$,
where the $b_m\ge0$ are densities on $(0,u_L]$ and $c_{kj,m}\ge0$, so that
$\alpha_{kj}=\sum_mc_{kj,m}$. Normalised indicators of the bins $[u_{m-1},u_m)$ give the
histogram kernels of the model-independent stochastic declustering of \citet{marsan2008}.
B-spline densities give smooth kernels. In both cases $\varphi_{kj}$ is a nonnegative
mixture. Adding the mixture component $m$ to the parent label therefore keeps every
M-step in closed form:
\[
  c_{kj,m}=\frac{N_{kj,m}}{\sum_{l}\int_0^{T-t^j_l}b_m(u)\,\dd u},\qquad
  N_{kj,m}=\sum_{i,l}\frac{c_{kj,m}b_m(t^k_i-t^j_l)}{\lambda_k(t^k_i)} .
\]
The denominator is the window correction of \cref{prop:est-mstep}(b), applied
component by component. A roughness penalty on $(c_{kj,m})_m$ can be added through the
one-step-late EM \citep{green1990} or a penalised Newton step, as in the nonparametric EM
of \citet{lewis2011}. \Citet{zhou2013b} estimate multivariate triggering kernels by a
related majorise--minimise scheme. There is no recursion for these kernels. The E-step
sums over pairs closer than the support $u_L$, at a cost of $O(n_{\mathrm{ev}} w)$, where $w$ is the mean
number of incidents in a window of length $u_L$. We choose $u_L$ below the knot spacing
of the background. This builds the separation of time scales of \cref{sec:est-mle} into
the design. Otherwise the long tail of a free kernel and a slow background compete for the
same variation.

\paragraph{Sparse branching matrices.} For many streams, sparsity can also be imposed by
penalisation. \Citet{hansen2015} use the least-squares contrast
$\sum_k\{-2\int_0^T\lambda_k\,\dN_k+\int_0^T\lambda_k^2\,\dd t\}$. Its expectation is
minimised by the true intensity, and it is quadratic in the coefficients of a linear
intensity. With weighted $\ell_1$ penalties whose data-driven weights come from
martingale Bernstein inequalities, they obtain a Lasso with oracle inequalities.
\Citet{zhou2013} combine $\ell_1$ and nuclear-norm penalties on $\Amat$ for networks that
are sparse and of low rank. Within EM, an $\ell_1$ penalty $\lambda^{\mathrm{pen}}\alpha_{kj}$
keeps the M-step in closed form, $\alpha_{kj}=N_{kj}/(G_{kj}+\lambda^{\mathrm{pen}})$. It
shrinks every entry but never sets one exactly to zero. The Bayesian priors of
\cref{sec:est-bayes} quantify uncertainty about the sparsity pattern and are our default.

\subsection{The boundary likelihood-ratio test}\label{app:est-lrt}

We use the notation of \cref{sec:est-lrt}: $\vartheta=(\alpha,\vartheta_2)$ with
$\alpha=\alpha_{kj}$ and $\vartheta_2\in\Theta_2\subset\mathbb R^d$. The parameter space is
$\Theta=[0,\infty)\times\Theta_2$, intersected with constraints that are inactive at
$\vartheta_0$, and the null set is $\Theta_0=\{0\}\times\Theta_2$. Let $\hat\vartheta$ and
$\tilde\vartheta$ maximise $\ell_T$ over $\Theta$ and $\Theta_0$, and let
$\mathrm{LR}_T=2\{\ell_T(\hat\vartheta)-\ell_T(\tilde\vartheta)\}$. Let
$U_T=T^{-1/2}\nabla\ell_T(\vartheta_0)$, where the derivative in $\alpha$ is one-sided, and
$J_T(\vartheta)=-T^{-1}\nabla^2\ell_T(\vartheta)$. The dimension $d$ is fixed as
$T\to\infty$. We assume the following:
\begin{enumerate}[({R}1)]
  \item \emph{Consistency.} $\vartheta_0=(0,\xi_0)$ with $\xi_0$ interior to $\Xi$, and
    $\hat\vartheta\to\vartheta_0$ and $\tilde\vartheta\to\vartheta_0$ in probability.
  \item \emph{Score.} $U_T\Rightarrow N(0,\bm I)$ with $\bm I$ positive definite.
  \item \emph{Information.} For every $\delta_T\downarrow0$,
    $\sup\{\|J_T(\vartheta)-\bm I\|:\vartheta\in\Theta,\ |\vartheta-\vartheta_0|\le\delta_T\}\to0$
    in probability.
\end{enumerate}
Partition $\bm I$ conformably with $(\alpha,\xi)$, and let
$I_{\alpha\alpha\cdot\xi}=I_{\alpha\alpha}-I_{\alpha\xi}I_{\xi\xi}^{-1}I_{\xi\alpha}$ be
the efficient information for $\alpha$.

\begin{proposition}[Boundary likelihood-ratio test; \citealp{chernoff1954,self1987}]\label{prop:lrt}
Under {(R1)}--{(R3)}:
\begin{enumerate}[(a)]
  \item $\mathrm{LR}_T=(Z_T^+)^2+o_p(1)$, where
    $Z_T=(U_{T,\alpha}-I_{\alpha\xi}I_{\xi\xi}^{-1}U_{T,\xi})/I_{\alpha\alpha\cdot\xi}^{1/2}$
    is the standardised efficient score. Hence
    \[
      \mathrm{LR}_T\Rightarrow\tfrac12\chi^2_0+\tfrac12\chi^2_1 .
    \]
    The asymptotic $p$-value of an observed value $\mathrm{LR}_T=x>0$ is
    $\tfrac12\Pr(\chi^2_1\ge x)$, and it is $1$ if $x=0$. The $5\%$ critical
    value is $2.71$, not $3.84$.
  \item $\sqrt T\,\hat\alpha\Rightarrow\max(0,W)$ with
    $W\sim N(0,I_{\alpha\alpha\cdot\xi}^{-1})$. Wald intervals for $\alpha_{kj}$ are
    therefore invalid near the boundary.
\end{enumerate}
\end{proposition}

\begin{proof}
\emph{Step 1: rate.} Since $\vartheta_0\in\Theta$, we have
$\ell_T(\hat\vartheta)\ge\ell_T(\vartheta_0)$. Near $\vartheta_0$ the set $\Theta$ is
convex, so the segment from $\vartheta_0$ to $\hat\vartheta$ lies in $\Theta$. Taylor's
theorem with integral remainder, applied along the segment, gives with
$h=\sqrt T(\hat\vartheta-\vartheta_0)$
\[
  0\le\ell_T(\hat\vartheta)-\ell_T(\vartheta_0)=h^{\top}U_T-\tfrac12h^{\top}\bar J_Th,
  \qquad \bar J_T=2\int_0^1(1-s)\,J_T\big(\vartheta_0+s(\hat\vartheta-\vartheta_0)\big)\,\dd s .
\]
By (R1) there are $\delta_T\downarrow0$ with $\Pr(|\hat\vartheta-\vartheta_0|\le\delta_T)\to1$,
and then (R3) gives $\bar J_T=\bm I+o_p(1)$. Hence
$\tfrac12\{\lambda_{\min}(\bm I)-o_p(1)\}|h|^2\le|h|\,|U_T|$, and $|h|=O_p(1)$ by (R2). The
same argument applies to $\tilde\vartheta$.

\emph{Step 2: local quadratic approximation.} The same expansion shows that, uniformly in
$h\in\sqrt T(\Theta-\vartheta_0)$ with $|h|\le c$,
\[
  \ell_T(\vartheta_0+h/\sqrt T)-\ell_T(\vartheta_0)=\mathcal Q_T(h)+o_p(1),\qquad
  \mathcal Q_T(h)=h^{\top}U_T-\tfrac12h^{\top}\bm Ih .
\]
Because $\xi_0$ is interior, the only constraint active near $\vartheta_0$ is
$\alpha\ge0$. For each fixed $c$ and large $T$, the local parameter set inside
$\{|h|\le c\}$ is therefore $\mathcal C\cap\{|h|\le c\}$, where
$\mathcal C=[0,\infty)\times\mathbb R^d$. The local null set is likewise
$\mathcal L\cap\{|h|\le c\}$, where $\mathcal L=\{0\}\times\mathbb R^d$. The maximisers of
$\ell_T$ (Step 1) and of $\mathcal Q_T$ over $\mathcal C$ and over $\mathcal L$ are all $O_p(1)$.
Hence
$2\{\ell_T(\hat\vartheta)-\ell_T(\vartheta_0)\}=2\sup_{\mathcal C}\mathcal Q_T+o_p(1)$, and
likewise with $\tilde\vartheta$ and $\mathcal L$. Let $Y_T=\bm I^{-1}U_T$ and
$\|v\|^2=v^{\top}\bm Iv$. Then $2\mathcal Q_T(h)=\|Y_T\|^2-\|Y_T-h\|^2$, so
\[
  \mathrm{LR}_T=\inf_{h\in\mathcal L}\|Y_T-h\|^2-\inf_{h\in\mathcal C}\|Y_T-h\|^2+o_p(1).
\]

\emph{Step 3: geometry.} Let $v=\bm I^{-1}\bm e_1$, where $\bm e_1$ is the
$\alpha$-coordinate. For $h\in\mathcal L$, $\langle h,v\rangle=h^{\top}\bm I\bm I^{-1}\bm e_1=h_1=0$,
so $v$ is orthogonal to $\mathcal L$ in the inner product of $\bm I$. Every $h$
decomposes as $h=h_{\mathcal L}+t\,v/\|v\|$, with $t=h_1/\|v\|$ and
$\|v\|^2=(\bm I^{-1})_{11}>0$. Hence $\mathcal C=\mathcal L\oplus\{t\,v/\|v\|:t\ge0\}$.
Write $Y_T=Y_{\mathcal L}+s_T\,v/\|v\|$. Then
$\inf_{\mathcal L}\|Y_T-h\|^2=s_T^2$ and $\inf_{\mathcal C}\|Y_T-h\|^2=\min(s_T,0)^2$, so
$\mathrm{LR}_T=(s_T^+)^2+o_p(1)$. By the partitioned inverse,
$(\bm I^{-1})_{11}=I_{\alpha\alpha\cdot\xi}^{-1}$ and
$(\bm I^{-1}U_T)_1=I_{\alpha\alpha\cdot\xi}^{-1}(U_{T,\alpha}-I_{\alpha\xi}I_{\xi\xi}^{-1}U_{T,\xi})$,
so $s_T=(\bm I^{-1}U_T)_1/\|v\|=Z_T$. By (R2), $Z_T\Rightarrow N(0,1)$. Hence
$(Z_T^+)^2$ is $0$ with probability tending to $\tfrac12$ and is otherwise asymptotically
$\chi^2_1$.

(b) The maximiser of $\mathcal Q_T$ over $\mathcal C$ is the $\bm I$-projection
$Y_{\mathcal L}+s_T^+v/\|v\|$. Its $\alpha$-coordinate is
$s_T^+\|v\|=\max\{0,(\bm I^{-1}U_T)_1\}$. The limit $\mathcal Q(h)=h^{\top}U-\tfrac12h^{\top}\bm Ih$ has a unique maximiser over the
closed convex cone $\mathcal C$, and the approximation of Step 2 is uniform on compacts. The
argmax continuous mapping theorem \citep[Thm~3.2.2]{vandervaart1996} therefore gives the
claim, with $W=(\bm I^{-1}U)_1\sim N(0,(\bm I^{-1})_{11})$.
\end{proof}

\subsection{Decays identified only under the alternative}\label{app:est-davies}

Suppose $\varphi_{kj}$ has its
own decay rate. Restrict it to a pre-registered interval $\mathcal B=[\beta_L,\beta_U]$ of
half-lives, and let $\mathrm{LR}_T(\beta)$ be the statistic with $\beta_{kj}=\beta$ fixed. For
each $\beta$, \cref{prop:lrt} gives $\mathrm{LR}_T(\beta)=(Z_T(\beta)^+)^2+o_p(1)$. Under
uniform versions of (R1)--(R3), $Z_T(\cdot)$ converges to a centred, unit-variance Gaussian
process $Z$, whose correlation function is that of the efficient scores. Then
$\sup_\beta\mathrm{LR}_T(\beta)\Rightarrow\sup_\beta(Z(\beta)^+)^2$
\citep{davies1977,davies1987}. This limit is stochastically larger than
$\tfrac12\chi^2_0+\tfrac12\chi^2_1$, because it dominates $(Z(\beta_0)^+)^2$ for every
$\beta_0\in\mathcal B$. Referring the sup-statistic to the chi-bar-square therefore inflates
the size. We use two valid references.
\begin{itemize}
  \item \emph{An upcrossing bound.} If $Z$ has continuously differentiable paths, Rice's
    formula for the expected number of upcrossings gives, for $c>0$,
    \[
      \Pr\Big(\sup_{\beta\in\mathcal B}Z(\beta)>c\Big)\le\Phi(-c)+\frac{\mathcal V}{2\pi}e^{-c^2/2},\qquad
      \mathcal V=\int_{\mathcal B}\{\Var Z'(\beta)\}^{1/2}\,\dd\beta .
    \]
    Since $\E|Z'(\beta)|=(2/\pi)^{1/2}\{\Var Z'(\beta)\}^{1/2}$, we can estimate
    $\mathcal V$ by $(\pi/2)^{1/2}$ times the total variation of the observed path on a
    grid $\beta_1<\dots<\beta_G$. With $u=\sup_\beta\mathrm{LR}_T(\beta)$ and
    $\widehat{\mathrm{TV}}=\sum_g|\hat Z(\beta_g)-\hat Z(\beta_{g-1})|$, this gives
    \[
      p_{\mathrm D}=\tfrac12\Pr(\chi^2_1>u)+\frac{\widehat{\mathrm{TV}}}{2\sqrt{2\pi}}\,e^{-u/2}.
    \]
    This is the one-sided form of Davies' bound, half of his two-sided bound. It is an
    upper bound only when $\mathcal V$ is known. The grid estimate of $\mathcal V$ is
    random and biased downward, because a coarse grid misses oscillations of the path. So
    $p_{\mathrm D}$ is an approximation, not a guaranteed upper bound, and a fine grid
    reduces the bias. For $\hat Z(\beta)$ we recommend the standardised efficient score at
    the null fit. It needs a single null fit and one recursion \eqref{eq:est-recursion} per
    grid point. It also follows the whole Gaussian path, whereas the signed root of
    $\mathrm{LR}_T(\beta)$ is flattened at zero wherever $\hat\alpha_{kj}=0$, which
    understates the total variation.
  \item \emph{Parametric bootstrap.} Simulate the full model at the null fit
    $\tilde\vartheta$ (\cref{app:computation}). Recompute the sup-statistic with refits,
    and report $p_{\mathrm B}=(1+\#\{\mathrm{LR}^{(b)}_{\sup}\ge\mathrm{LR}_{\sup}\})/(B+1)$ for
    $B$ replicates, where $\mathrm{LR}_{\sup}=\sup_\beta\mathrm{LR}_T(\beta)$.
\end{itemize}

\subsection{Time rescaling}\label{app:est-gof}

\begin{proposition}[Time rescaling; \citealp{papangelou1972,brown2002}]\label{prop:rescaling}
Let $N$ be a simple point process on $[0,\infty)$ with intensity $\lambda>0$ with respect
to a right-continuous filtration $(\Hist_t)$. Suppose $\Lam(t)=\int_0^t\lambda(s)\,\dd s$ is
finite for every $t$ and $\Lam(t)\to\infty$ almost surely. Then the points
$\Lam(t_1)<\Lam(t_2)<\cdots$ form a homogeneous Poisson process with unit rate. In
particular, the increments $z_i=\Lam(t_i)-\Lam(t_{i-1})$, with $t_0=0$, are i.i.d.\
$\mathrm{Exp}(1)$, and the $u_i=1-e^{-z_i}$ are i.i.d.\ uniform on $(0,1)$. If the
multivariate model is correct, this holds for each stream $N_k$ with its compensator
$\Lam_k$ relative to the joint history. Here $\lambda_k\ge\mu_k>0$, so the positivity
condition holds.
\end{proposition}

\begin{proof}
Since $\lambda>0$, $\Lam$ is continuous and strictly increasing from $0$ to $\infty$, so it
has a continuous inverse $\Lam^{-1}$. Each $\Lam^{-1}(u)$ is a stopping time, because
$\{\Lam^{-1}(u)\le t\}=\{\Lam(t)\ge u\}\in\Hist_t$. Let $\tilde N(u)=N(\Lam^{-1}(u))$ and
$\tilde\Hist_u=\Hist_{\Lam^{-1}(u)}$. Then $\tilde N(u)=\#\{i:\Lam(t_i)\le u\}$, and
$\tilde N$ is simple because $\Lam$ is injective.

We show that $\tilde N(u)-u$ is a $\tilde\Hist$-martingale. Fix $v$. The stopped process
$N(t\wedge\Lam^{-1}(v))-\Lam(t\wedge\Lam^{-1}(v))$ is a local martingale whose compensator
part is bounded by $v$. By localisation and monotone convergence,
$\E N(\Lam^{-1}(v))=v$. The stopped process is therefore dominated by the integrable
$N(\Lam^{-1}(v))+v$, so it is a uniformly integrable martingale. Optional sampling at
$\Lam^{-1}(u)\le\Lam^{-1}(v)$ gives
$\E\{\tilde N(v)-\tilde N(u)\mid\tilde\Hist_u\}=v-u$.

A simple counting process whose compensator is the identity is a unit-rate Poisson process.
This is Watanabe's characterisation \citep[Ch.~II]{bremaud1981}. For the multivariate
model, apply the argument to each $N_k$ with the filtration $\Hist$. The positivity
assumption only simplifies the proof: the result holds for any continuous compensator
\citep[Sec.~7.4]{daley2003}.
\end{proof}

On a finite window, the rescaled points of stream $k$ are those of a unit Poisson process
on $[0,\Lam_k(T)]$, where $\Lam_k(T)$ is a stopping time for the time-changed filtration.
The $n_k$ increments are therefore not exactly a fixed-size i.i.d.\ sample, and the final
censored gap is ignored. The discrepancy vanishes as $\Lam_k(T)\to\infty$.

\paragraph{Estimated parameters.} In the i.i.d.\ setting with maximum likelihood
estimates, the direction of the error from plugging in $\hat\vartheta$ is known. The
estimated empirical process converges to a Gaussian process whose covariance is that of the
Brownian bridge minus a positive-semidefinite kernel \citep{durbin1973}. By Anderson's
inequality \citep{anderson1955}, a centred Gaussian law with smaller covariance puts more
mass on the symmetric convex set $\{\sup|z|\le c\}$. The plug-in statistic is therefore
asymptotically stochastically smaller, and the test is conservative, as in
\citet{lilliefors1967}. We know of no proof for rescaled point processes, but the same
mechanism operates when $\hat\vartheta$ is the maximum likelihood estimator. The argument
can fail for penalised, least-squares or posterior-mean plug-ins, whose covariance
correction need not be positive semidefinite. Hence the parametric bootstrap of
\cref{sec:est-gof}.

\subsection{Proofs for Section~\ref{sec:estimation}}\label{app:est-proofs}

\begin{proof}[Proof of \cref{prop:est-mstep}]
$Q$ separates into a background term for each $k$, a shock term for each $k$, and a term
for each pair $(k,j)$.

(a) The background term is $\sum_ip^{k,0}_i\log\mu_k(t^k_i)-\int_0^T\mu_k$. For constant
$\mu_k$ it equals $N^0_k\log\mu_k-T\mu_k$. For the log-linear form, the first part is linear
in $\theta_k$. The integral $\int e^{B^{\top}\theta}$ is strictly convex when the basis is
linearly independent, and the penalty is convex.

(b) The pair term is $N_{kj}\log\alpha_{kj}-\alpha_{kj}\mathcal G_{kj}$, plus terms free of
$\alpha_{kj}$. It is maximised at $N_{kj}/\mathcal G_{kj}$. The shock term in $\beta_k$ is
$(\sum_eN_{ke})\log\beta_k-\beta_k\sum_eS_e\Psi^T_k(e)$.

(c) For $g_{kj}(u)=\beta e^{-\beta u}$ the pair term is
$N\log\alpha+N\log\beta-\beta W-\alpha\beta\Upsilon(\beta)$, with $N=N_{kj}$, $W=W_{kj}$ and
$\Upsilon(\beta)=\sum_l\int_0^{r_l}e^{-\beta u}\,\dd u$; note that $\mathcal G_{kj}=\beta\Upsilon(\beta)$.
Profiling $\alpha=N/\{\beta\Upsilon(\beta)\}$ gives $-\beta W-N\log\Upsilon(\beta)$ plus a constant.
Now $\Upsilon$ is the Laplace transform of the measure $\nu(\dd u)=\sum_l\ind{0\le u\le r_l}\dd u$.
Laplace transforms are log-convex, and here strictly so, because $\nu$ is not a point mass:
$(\log\Upsilon)''(\beta)=\Var_\beta(u)>0$, where $\Var_\beta$ is the variance under
$e^{-\beta u}\nu(\dd u)/\Upsilon(\beta)$. Hence the profiled objective is strictly concave, and
its derivative is $-W+N\,m_j(\beta)$, since $m_j=-\Upsilon'/\Upsilon$ is the mean of $u$ under the tilted
measure. Moreover $m_j'(\beta)=-\Var_\beta(u)<0$. As $\beta\to\infty$ the tilted measure
concentrates at $0$. As $\beta\downarrow0$ it tends to $\nu/\nu(\mathbb R)$, whose mean is
$\sum_lr_l^2/(2\sum_lr_l)$. A shared decay sums strictly concave terms. The response rate
is the same calculation with $\nu=\sum_eS_e\ind{0\le u\le T-\tau_e}\dd u$ for
$\tau_e\ge0$.

(d) This is the usual EM argument \citep{dempster1977}. Write
$\ell_T(\vartheta)=Q(\vartheta)-H(\vartheta)$, where
$H(\vartheta)=\E\{\log p_\vartheta(Z\mid\text{data})\mid\text{data},\vartheta^{(r)}\}$. By
Gibbs' inequality, $H(\vartheta)\le H(\vartheta^{(r)})$. The penalty is a deterministic
function of $\vartheta$, so it passes through unchanged.
\end{proof}

\begin{proof}[Proof of \cref{lem:est-trunc}]
Write $\bm A_x=\Amat^{(-kj)}+x\bm e_k\bm e_j^{\top}$ and $\rho(x)=\spr(\bm A_x)$, so
$\rho(0)=\spr(\Amat^{(-kj)})<1$. By Perron--Frobenius
monotonicity, $\rho$ is nondecreasing and continuous on $[0,\infty)$. Let
$x^*=\sup\{x:\rho(x)<1\}$. If $\rho(x)<1$, the eigenvalues $\lambda_i$ of $\bm A_x$ satisfy
$\det(\Imat-\bm A_x)=\prod_i(1-\lambda_i)>0$: complex eigenvalues come in conjugate pairs,
and real ones are below $1$. By the matrix determinant lemma,
\[
  \det(\Imat-\bm A_x)=\det(\Imat-\Amat^{(-kj)})\,\big(1-x\,\bm e_j^{\top}\Mmat^{(-kj)}\bm e_k\big),
\]
which is affine in $x$ and positive at $x=0$. Suppose first that $x^*<\infty$. Then
$\rho(x^*)=1$ by continuity, and $1$ is an eigenvalue of the nonnegative matrix
$\bm A_{x^*}$ by Perron--Frobenius. So the determinant vanishes at $x^*$, and $x^*$ is its
unique root $1/\Mmat^{(-kj)}_{jk}$. Conversely, $\Mmat^{(-kj)}=\sum_{m\ge0}(\Amat^{(-kj)})^m\ge\zero$.
If $\Mmat^{(-kj)}_{jk}>0$, then $x^*<\infty$: otherwise the determinant would be positive
for every $x$, yet it vanishes at $x=1/\Mmat^{(-kj)}_{jk}$. If $\Mmat^{(-kj)}_{jk}=0$, the
determinant never vanishes, so $x^*=\infty$.
\end{proof}

\begin{proof}[Proof of \cref{prop:est-binning}]
(a) Write $\pi(z)=1-\sum_\ell\pi_\ell z^\ell$ and $f_X$ for the spectral density of $X$.
The Wold representation gives $f_X(\omega)=\sigma^2_\Delta/\{2\pi|\pi(e^{-i\omega})|^2\}$
\citep[Ch.~5]{brockwell1991}. At $\omega=0$,
$2\pi f_X(0)=\sum_h\Cov(X_0,X_h)=\lim_m m^{-1}\Var N(0,m\Delta]=F_\infty\bar\Lambda\Delta$,
and $F_\infty=(1-n)^{-2}$ by \cref{prop:th-spectrum}. The function $\pi$ is continuous
on the closed unit disc. It has no zeros there: inside, $\pi\psi=1$ for the Wold transfer
function $\psi$, and on the circle $|\pi|^2=\sigma^2_\Delta/(2\pi f_X)>0$ because $f_X$ is
bounded. It is real on $[0,1]$ with $\pi(0)=1$, so $\pi(1)>0$. Hence
\[
  1-n_\Delta=\pi(1)=\frac{\sigma_\Delta}{\{2\pi f_X(0)\}^{1/2}}=(1-n)\,\varpi_\Delta^{1/2}.
\]
For $\varpi_\Delta\ge1$, note that the predictor is a function of $\Hist_{(t-1)\Delta}$,
so $\sigma^2_\Delta\ge\E\Var(X_t\mid\Hist_{(t-1)\Delta})$. Given the past, the future of a
Hawkes process is again a Poisson cluster process. Its cluster seeds form a Poisson
process with the deterministic intensity
$\mu+\int_{(-\infty,(t-1)\Delta]}\varphi(\cdot-s)\,\dN(s)$, and each seed carries an
independent cluster \citep{hawkes1974}. So $X_t$ is conditionally a Poisson sum
$\sum_cS_c$ of within-bin cluster sizes $S_c\ge1$. Its conditional variance is
$\E[\sum_cS_c^2\mid\Hist]\ge\E[\sum_cS_c\mid\Hist]=\E[X_t\mid\Hist]$. Taking expectations
gives $\sigma^2_\Delta\ge\bar\Lambda\Delta$.

(b) The covariance density is $c(u)=Ke^{-\gamma|u|}$ by \eqref{eq:th-fH}. Integrating it
over pairs of bins gives $V$ and $\Cov(X_t,X_{t+\ell})=K(1-e^{-\gamma\Delta})^2e^{-\gamma\Delta(\ell-1)}/\gamma^2=C\rho^\ell$.
Hence $Y_t=X_t-\rho X_{t-1}$ has $\Var Y_t=v_0$, $\Cov(Y_t,Y_{t-1})=v_1$, and zero
autocovariance at lags $\ge2$. So $Y$ is an MA(1) process, $Y_t=\varepsilon_t-\zeta\varepsilon_{t-1}$.
Binning and aliasing map spectra to spectra monotonically, and the Bartlett spectrum
\eqref{eq:th-fH} dominates the Poisson floor $\bar\Lambda/(2\pi)$. Hence
$f_X\ge\bar\Lambda\Delta/(2\pi)>0$, the binned spectrum of a Poisson process. So
$f_Y(\omega)=|1-\rho e^{-i\omega}|^2f_X(\omega)=v_0+2v_1\cos\omega>0$, and the invertible
solution has $|\zeta|<1$. Matching $\sigma^2(1+\zeta^2)=v_0$ and $-\zeta\sigma^2=v_1$ shows
that $\sigma^2$ is the larger root of $s^2-v_0s+v_1^2=0$. Inverting
$(1-\rho z)/(1-\zeta z)=\pi(z)$ gives $\pi_\ell$ and $n_\Delta$. For the expansion, put
$u=V-C=\bar\Lambda\Delta+O(\Delta^3)$ and note $C=K\Delta^2+O(\Delta^4)$. Then
$v_0=u(1+\rho^2)+C(1-\rho^2)=2\bar\Lambda\Delta-2\gamma\bar\Lambda\Delta^2+O(\Delta^3)$, and
\[
  v_0^2-4v_1^2=(1-\rho^2)\{u^2(1-\rho^2)+2uC(1+\rho^2)+C^2(1-\rho^2)\}
  =4\gamma\bar\Lambda(\gamma\bar\Lambda+2K)\Delta^4+O(\Delta^5).
\]
Since $\gamma\bar\Lambda+2K=\bar\Lambda\beta/(1-n)$, the square root is
$2\bar\Lambda\beta\Delta^2+O(\Delta^3)$. So
$\sigma^2_\Delta=\bar\Lambda\Delta\{1+(\beta-\gamma)\Delta\}+O(\Delta^3)$, with
$\beta-\gamma=n\beta$, and (a) gives the expansion of $n_\Delta$.
\end{proof}

\begin{proof}[Proof of \cref{lem:est-placebo}]
Given the conditioning variables, $\Gamma^{(b)}=\Gamma(C_b;\text{data})$ applies one fixed
function to exchangeable arguments, so $(\Gamma^{(0)},\dots,\Gamma^{(B)})$ is exchangeable.
Let $R_b=\#\{b'\in\{0,\dots,B\}:\Gamma^{(b')}\ge\Gamma^{(b)}\}$, so that $p=R_0/(B+1)$.
Order the $B+1$ values decreasingly, breaking ties arbitrarily. The position of
$\Gamma^{(b)}$ is at most $R_b$, and distinct indices have distinct positions, so at most
$m$ indices satisfy $R_b\le m$. By exchangeability,
$\Pr(R_0\le m)=(B+1)^{-1}\sum_b\Pr(R_b\le m)=(B+1)^{-1}\E\#\{b:R_b\le m\}\le m/(B+1)$. Take
$m=\lfloor a(B+1)\rfloor$, and integrate over the conditioning variables.
\end{proof}

\section{Supplement to Section~\ref{sec:forecasting}}\label{app:fc-supp}

This appendix collects material behind \cref{sec:forecasting}: exact predictive
distributions, an imputation algorithm for alerting from delayed data, the link between
cost-weighted losses and the Brier score, and a monitor for regime shifts. Proofs are in
\cref{app:attrproofs}. None of this material is evaluated in the simulation study.

\subsection{Exact predictive distributions}\label{app:fc-pgf}

\begin{remark}[Exact predictive distributions]\label{rem:fc-pgf}
The predictive law can also be computed without simulation. Fix a target stream $k$
and the window $W=(t,t+\Delta_{\mathrm a}]$. For a cluster rooted at a stream-$j$ incident at time
$s$, let $G_j(z;s)$ be the probability generating function of the number of its members
in stream $k$ and in $W$, root included. The branching structure gives
\begin{equation}\label{eq:fc-pgf-cluster}
  G_j(z;s)=z^{\ind{j=k,\ s\in W}}\exp\Bigl\{\sum_{i=1}^{M}\int_s^{t+\Delta_{\mathrm a}}\varphi_{ij}(r-s)\bigl(G_i(z;r)-1\bigr)\,\dd r\Bigr\},
\end{equation}
and the Poisson immigrants of \cref{lem:fc-future} give
\begin{equation}\label{eq:fc-pgf}
  \E\bigl[z^{N_k(W)}\mid\Hist_t\bigr]
  =\exp\Bigl\{\sum_{j=1}^{M}\int_t^{t+\Delta_{\mathrm a}}\nu_{t,j}(s)\bigl(G_j(z;s)-1\bigr)\,\dd s\Bigr\}.
\end{equation}
For exponential kernels, the inner integrals in \eqref{eq:fc-pgf-cluster} satisfy $M^2$
ordinary differential equations, solved backwards in $s$. In the univariate case,
$I(s)=\int_s^{t+\Delta_{\mathrm a}}\varphi(r-s)\{G(z;r)-1\}\,\dd r$ solves
\[
I'=\beta I-n\beta\,(ze^{I}-1),\qquad I(t+\Delta_{\mathrm a})=0,
\]
and $G=ze^{I}$ on $W$. Evaluating \eqref{eq:fc-pgf} at the $L$th roots of unity and
applying the discrete Fourier transform gives the predictive probability mass function.
The aliasing error in each probability is at most the predictive probability of $L$ or
more incidents.

Because $\bm\nu_t$ is affine in $\bm\xi(t)$, the logarithm of \eqref{eq:fc-pgf} is
affine in $\bm\xi(t)$. The predictive law is therefore the law at $\bm\xi=\zero$
convolved with independent compound Poisson laws whose rates are proportional to the
coordinates of $\bm\xi(t)$. This is the analytic form of the coupling in the proof of
\cref{thm:alert}.
\end{remark}

\subsection{Alerting from delayed data by imputation}\label{app:fc-sis}

We use the notation of \cref{sec:fc-nowcast}. The plug-in rule that compares a nowcast
state with the threshold curve of \cref{thm:alert}(d) is not the Bayes rule for delayed
data, because $G_t$ is nonlinear. The following lemma and \cref{alg:alert} approximate
the predictive probability given the available data instead.

\begin{lemma}[Importance weights for imputation]\label{lem:fc-is}
Assume \cref{ass:th-model,ass:fc-report} with known parameters, and condition on the
history up to $t-L$. Let $O$ be the incidents in $(t-L,t]$ that are available by $t$.
Simulate the remaining incidents $U$ in $(t-L,t]$ by thinning with intensity
$\{1-p_kF_k(t-s)\}\lambda_k(s)$, where $\lambda$ is computed from the history up to
$t-L$, from $O$ and from the points of $U$ simulated so far. Then the conditional law of
$U$ given $O$ has density, with respect to this proposal, proportional to
\begin{equation}\label{eq:fc-isweight}
  w(U)=\prod_{i\in O}\lambda_{k_i}(s_i)\,
  \exp\Bigl\{-\sum_{k=1}^{M}\int_{t-L}^{t}p_kF_k(t-s)\,\lambda_k(s)\,\dd s\Bigr\}.
\end{equation}
Consequently, if $U_1,\dots,U_J$ are independent proposals and $\bm\xi_\ell(t)$ is the
state of the completed history, then $\sum_\ell w(U_\ell)G_t(\bm\xi_\ell(t))/\sum_\ell w(U_\ell)$
converges almost surely, as $J\to\infty$, to $\Pr(Y_t=1\mid O,\ \text{history up to }t-L)$.
\end{lemma}

In practice the history up to $t-L$ is itself incomplete. \Cref{alg:alert} replaces it
by the weighted carried-over excitation \eqref{eq:fc-ht} with $\delta=L$. The window $L$
is chosen so that $F_k(L)\approx1$ and the excitation carried over from before $t-L$ has
largely decayed by $t$. The only approximations, beyond Monte Carlo error, are this
substitution and the parameter draws. When $p_k<1$, unrecorded incidents never become
available. A fraction $1-p_k$ of the intensity is then imputed throughout the window,
not only near $t$.

\begin{algorithm}[t]
\caption{Real-time alerting with imputation of unreported incidents}\label{alg:alert}
\begin{algorithmic}[1]
\Require incidents available at $t$; parameter draws $\vartheta^{(1)},\dots,\vartheta^{(R)}$;
  recording probabilities $p_k$ and delay distributions $F_k$; imputation window $L$
  with $F_k(L)\approx1$; window $\Delta_{\mathrm a}$; surge set $\mathcal U$; threshold
  $\pi^*=c_{\mathrm{FA}}/(c_{\mathrm{FA}}+c_{\mathrm{miss}})$; particles per draw $J$
\Ensure alert $a_t$ and estimate $\hat\pi_t$ of $\Pr(Y_t=1\mid\mathcal G_t)$
\For{$r=1,\dots,R$}
  \State $\hat{\bm h}^{(r)}\gets$ \eqref{eq:fc-ht} with $\delta=L$, under $\vartheta^{(r)}$
    \Comment{state at $t-L$}
  \For{$\ell=1,\dots,J$}
    \State impute the unavailable incidents in $(t-L,t]$ by thinning with intensity
      $\{1-p_kF_k(t-s)\}\lambda_k(s)$, inserting the available ones
    \State $w_{r\ell}\gets$ \eqref{eq:fc-isweight} for the completed path
      \Comment{\cref{lem:fc-is}}
    \State $\pi_{r\ell}\gets G_t(\bm\xi_{r\ell}(t))$ for the completed state
      \Comment{simulation or \cref{rem:fc-pgf}}
  \EndFor
  \State $\hat\pi^{(r)}\gets\sum_\ell w_{r\ell}\pi_{r\ell}\big/\sum_\ell w_{r\ell}$
\EndFor
\State $\hat\pi_t\gets R^{-1}\sum_r\hat\pi^{(r)}$
\State \Return $a_t=\ind{\hat\pi_t\ge\pi^*}$ and $\hat\pi_t$
\end{algorithmic}
\end{algorithm}

Several points about \cref{alg:alert} follow.
\begin{itemize}
  \item \emph{Normalising the weights.} Weights are normalised within each parameter
  draw. The draws already condition on all the data, so weighting across draws would
  count the recent data twice.
  \item \emph{Effective sample size.} The effective sample size
  $(\sum_\ell w_{r\ell})^2/\sum_\ell w_{r\ell}^2$ should be monitored. If it is small,
  $J$ should be increased or $L$ shortened.
  \item \emph{Monte Carlo error.} The Monte Carlo standard error of $\hat\pi_t$ is
  estimated from the spread of the $R$ per-draw estimates, or by the delta method for
  ratio estimators when $R=1$. If $|\hat\pi_t-\pi^*|$ is within two standard errors,
  the simulation size should be increased before deciding.
  \item \emph{Parameter uncertainty.} With posterior draws of $\vartheta$, $p_k$ and $F_k$,
  $\hat\pi_t$ estimates the posterior predictive probability. This is $\pi_t$ in
  \cref{thm:alert}(a) when the information includes parameter uncertainty.
  \item \emph{Calibration in practice.} The algorithm is calibrated only if the model
  is. Before operational use, the PIT (\cref{lem:fc-pit}) and a reliability diagram of
  $\hat\pi_t$ should be checked on held-out data. If they show miscalibration,
  $\hat\pi_t$ should be recalibrated, for instance by isotonic regression
  (\cref{rem:fc-brier}), before the threshold $\pi^*$ is applied.
  \item \emph{Plug-in shortcuts.} Comparing a nowcast state with the threshold curve of
  \cref{thm:alert}(d) is fast but is not the Bayes rule for delayed data. If it is
  used, it should be recalibrated in the same way.
  \item \emph{Reuse.} With $G_t$ replaced by simulated future counts, the same weighted
  particles give the predictive distributions used for scoring in
  \cref{sec:fc-eval}.
\end{itemize}

\subsection{Cost-weighted losses and the Brier score}\label{app:fc-brier}

\begin{remark}[Cost-weighted losses and the Brier score]\label{rem:fc-brier}
Normalise the costs so that $c_{\mathrm{FA}}=c$ and $c_{\mathrm{miss}}=1-c$. Then
$\pi^*=c$, and the loss of the rule $\ind{\pi_t\ge c}$ is
$S_c(\pi_t,Y_t)=c\,(1-Y_t)\ind{\pi_t\ge c}+(1-c)\,Y_t\ind{\pi_t<c}$. For all $\pi_t$
and $Y_t$,
\[
\int_0^1S_c(\pi_t,Y_t)\,\dd c=\tfrac12(\pi_t-Y_t)^2 .
\]
The Brier score is therefore twice the average realised cost when the normalised cost
$c=\pi^*$ is drawn uniformly from $(0,1)$. Every proper scoring rule for binary events is a mixture
of such elementary losses \citep{schervish1989}. When partners cannot agree on one cost
ratio, the Brier score summarises performance across normalised costs. When they can, the
realised cost at $\pi^*$ is the relevant measure.

\Cref{thm:alert}(a) presumes that $\pi_t$ is the true conditional probability. A
forecast that is not calibrated should therefore be recalibrated before the threshold
$\pi^*$ is applied, for instance by isotonic regression of $Y_t$ on $\pi_t$ over the
training period.
\end{remark}

\subsection{Detecting regime shifts}\label{sec:fc-cusum}

Alerts concern the next $\Delta_{\mathrm a}$ days with the parameters held fixed. A different
question is whether the parameters themselves have changed, through a sustained rise
either in the background or in contagion.

\begin{remark}[A CUSUM monitor]\label{rem:fc-cusum}
For a pre-change model $\vartheta_0$ and an alternative $\vartheta_1$, let
\[
\ell(t)=\sum_{k=1}^{M}\Bigl[\sum_{i:\,t^k_i\le t}\log\frac{\lambda^{\vartheta_1}_k(t^k_i)}{\lambda^{\vartheta_0}_k(t^k_i)}
-\int_0^t\bigl\{\lambda^{\vartheta_1}_k(s)-\lambda^{\vartheta_0}_k(s)\bigr\}\,\dd s\Bigr]
\]
be the log-likelihood ratio of the data on $[0,t]$, with both intensities computed from
the same observed history. Page's CUSUM statistic
$\mathcal W(t)=\ell(t)-\min_{s\le t}\ell(s)$ signals at the first $t$ with
$\mathcal W(t)\ge h$ \citep{page1954}. Two detectors can run in parallel:
\begin{itemize}
  \item one whose alternative is a shift in the background, $\bm\mu\mapsto(1+\epsilon_\mu)\bm\mu$
  with $\Amat$ fixed;
  \item one whose alternative is a shift in contagion,
  $\Amat\mapsto(1+\epsilon_A)\Amat$ with $\bm\mu$ fixed.
\end{itemize}
Comparing them shows which kind of change better explains a sustained departure. The
threshold $h$ is calibrated by simulation under $\vartheta_0$ to a target average run
length to false alarm. The optimality theory of the CUSUM is for independent
observations; here it serves as a monitoring statistic.

Two cautions apply. First, the detectors inherit the identification limits of
\cref{sec:identification}. Over short stretches, a shift in the background and a shift
in contagion can produce similar data, so the comparison is a diagnostic, not a test
that separates the two. Second, the statistics must be computed on settled data or with
the nowcast of \cref{prop:nowcast}, because reporting delays mimic a recent decline.
\end{remark}
